\chapter{Complex Analysis for Computing}

\begin{center}\itshape
The shortest path between two truths in the real domain often passes through the complex domain.\\
\upshape --- Paul Painlev\'e (1900), made famous by Jacques Hadamard
\end{center}

A function of a complex variable that has a derivative in the complex sense is astonishingly rigid. It is automatically
infinitely differentiable, it is the sum of its Taylor series on every disc where it is defined, and its values inside a
circle are determined by its values on the circle. Most important for us, its integral around a closed contour is
computed by looking at a few special points inside, the residues. This chapter develops just enough of the theory to use
it as a machine: real integrals, infinite sums, the extension of $\Gamma$ to the complex plane (which Chapter~10 needs for
the zeta function), and the counting of zeros.

On the way we keep several promises: the radius of convergence of a power series is the distance to the nearest
singularity (announced in Chapters~1 and~3), the partial fractions of $\pi\cot\pi z$ and the sixth proof of
$\zeta(2)=\frac{\pi^2}6$ come from residues (Chapter~3), and the transform of $\frac1{1+x^2}$ is recomputed (Chapter~8).

\begin{think}
The function $\frac1{1+x^2}$ is infinitely differentiable on the whole real line and has no visible problem anywhere. Yet its
Taylor series at $0$, $1-x^2+x^4-\cdots$, converges only for $|x|<1$. What happens at distance $1$ from $0$?
\end{think}

\section{Holomorphic functions and power series}

Throughout, $U$ denotes an open subset of $\C$, $D(a,r)=\{z:|z-a|<r\}$ the open disc and $\overline D(a,r)$ the closed disc. We
identify $\C$ with $\R^2$, so continuity, open sets and compactness have their usual meaning.

\begin{definition}
A function $f:U\to\C$ is \emph{complex differentiable} at $z_0\in U$ if
\[
f'(z_0)=\lim_{h\to0}\frac{f(z_0+h)-f(z_0)}h
\]
exists, $h$ tending to $0$ in $\C$. It is \emph{holomorphic} on $U$ if it is complex differentiable at every point of $U$. A
function holomorphic on all of $\C$ is \emph{entire}.
\end{definition}

The rules for sums, products, quotients and compositions have the same proofs as for real functions, so polynomials are
entire and rational functions are holomorphic away from the zeros of their denominators. The novelty is that $h$ may approach
$0$ from every direction, and the limit must not depend on the direction.

\begin{proposition}[Cauchy--Riemann equations]\label{prop:cr}
Write $f=u+iv$ with $u,v$ real, as functions of $(x,y)$ with $z=x+iy$. If $f$ is complex differentiable at $z_0$, then the partial
derivatives of $u$ and $v$ exist at $z_0$ and
\[
u_x=v_y,\qquad u_y=-v_x,\qquad f'(z_0)=u_x+iv_x.
\]
\end{proposition}
\begin{proof}
Let $h$ tend to $0$ along the real axis: $\frac{f(z_0+h)-f(z_0)}h\to u_x+iv_x$. Along the imaginary axis, $h=ik$ with $k$ real:
$\frac{f(z_0+ik)-f(z_0)}{ik}\to\frac{u_y+iv_y}i=v_y-iu_y$. Both limits equal $f'(z_0)$; compare real and imaginary parts.
\end{proof}

So $\overline z$ is nowhere complex differentiable ($u=x$, $v=-y$, $u_x=1\ne-1=v_y$), nor is $|z|^2$ except at $0$, although both are
smooth as functions of two real variables. Conversely, a function whose $u$ and $v$ have continuous partial derivatives
satisfying the Cauchy--Riemann equations is holomorphic; we will not need this.

\subsection*{Power series}

\begin{lemma}[Abel]\label{lem:abel}
Let $(c_n)$ be complex numbers and set
\[
R=\sup\{r\ge0:(|c_n|r^n)_n\text{ is bounded}\}\in[0,\infty].
\]
Then $\sum c_nz^n$ converges
absolutely for $|z|<R$, uniformly on every closed disc $\overline D(0,r)$ with $r<R$, and its terms are unbounded for $|z|>R$.
\end{lemma}
\begin{proof}
If $|z|>R$, the terms are unbounded by definition of $R$. If $r<\rho<R$, there is $M$ with $|c_n|\rho^n\le M$, so for $|z|\le r$,
$|c_nz^n|\le M(r/\rho)^n$, a convergent geometric series independent of $z$; apply the Weierstrass test (Lemma~\ref{lem:uniform}, whose
proof works on any set).
\end{proof}

$R$ is the \emph{radius of convergence}. Inside $D(0,R)$ the series converges; outside, it diverges since its terms do not
tend to $0$; on the circle anything can happen.

\begin{theorem}\label{thm:powerseries}
Let $f(z)=\sum_{n\ge0}c_n(z-a)^n$ have radius of convergence $R>0$. Then $f$ is holomorphic on $D(a,R)$ and
\[
f'(z)=\sum_{n\ge1}nc_n(z-a)^{n-1},
\]
a series with the same radius. Consequently $f$ has derivatives of all orders on $D(a,R)$, and $c_n=\frac{f^{(n)}(a)}{n!}$.
\end{theorem}
\begin{proof}
Take $a=0$. If $r<\rho<R$ and $|c_n|\rho^n\le M$, then $n|c_n|r^{n-1}\le\frac M\rho\,n(r/\rho)^{n-1}$, which is summable; so the derived series
converges absolutely on $\overline D(0,r)$ and its radius is at least $R$ (it is at most $R$ since $|c_nz^n|\le|z|\cdot|nc_nz^{n-1}|$ for
$n\ge1$).

Fix $z$ with $|z|<r<R$. For $|w|\le r$ let $q_n(w)=\sum_{k=0}^{n-1}w^kz^{n-1-k}$, so that $q_n(w)(w-z)=w^n-z^n$, $q_n(z)=nz^{n-1}$
and $|q_n(w)|\le nr^{n-1}$. The series $Q(w)=\sum_nc_nq_n(w)$ converges uniformly on $\overline D(0,r)$ by the bound just proved, so
$Q$ is continuous there. For $w\ne z$, $Q(w)=\frac{f(w)-f(z)}{w-z}$; letting $w\to z$ gives $f'(z)=Q(z)=\sum nc_nz^{n-1}$.

Applying the result to $f'$, then to $f''$, and so on, $f^{(k)}(z)=\sum_{n\ge k}n(n-1)\cdots(n-k+1)c_nz^{n-k}$, and $z=0$ gives
$f^{(k)}(0)=k!\,c_k$.
\end{proof}

\begin{example}[exponential and logarithm]
The series $\exp z=\sum\frac{z^n}{n!}$ has infinite radius, so $\exp$ is entire and $\exp'=\exp$. From $e^{a+b}=e^ae^b$ and
Proposition~2.1 of Volume~0, $e^{x+iy}=e^x(\cos y+i\sin y)$: $\exp$ never vanishes, and $e^z=1$ exactly when $z\in2\pi i\Z$.
Likewise $\cos z=\frac{e^{iz}+e^{-iz}}2$ and $\sin z=\frac{e^{iz}-e^{-iz}}{2i}$ are entire, and $\sin z=0$ means $e^{2iz}=1$, that is
$z\in\pi\Z$: the sine has no zeros off the real axis, so $\cot z=\frac{\cos z}{\sin z}$ is holomorphic on $\C\setminus\pi\Z$.

On the \emph{slit plane} $\C\setminus\left]-\infty,0\right]$, every $z$ is $re^{i\theta}$ with $r>0$ and a unique $\theta\in\left]-\pi,\pi\right[$, and
$\frac{\sin\theta}{1+\cos\theta}=\tan\frac\theta2$ gives the explicit continuous formula $\theta=2\arctan\frac{\operatorname{Im}z}{|z|+\operatorname{Re}z}$. The
\emph{principal logarithm} $\operatorname{Log}z=\log r+i\theta$ is thus continuous and inverts $\exp$. It is holomorphic with derivative
$\frac1z$: writing $\ell=\operatorname{Log}z$ and $\ell_w=\operatorname{Log}w$, $\frac{\ell_w-\ell}{w-z}=\frac{\ell_w-\ell}{e^{\ell_w}-e^\ell}\to\frac1{e^\ell}=\frac1z$ as $w\to z$,
since $\ell_w\to\ell$. Powers are defined by $z^s=e^{s\operatorname{Log}z}$ on the slit plane.
\end{example}

\section{Path integrals and Cauchy's theorem}

A \emph{path} is a continuous, piecewise $C^1$ map $\gamma:[\alpha,\beta]\to\C$; it is \emph{closed} if $\gamma(\alpha)=\gamma(\beta)$, and $\gamma^*$ denotes its
image. For $f$ continuous on $\gamma^*$,
\[
\begin{gathered}\int_\gamma f(z)\dd z=\int_\alpha^\beta f(\gamma(t))\gamma'(t)\dd t,\qq L(\gamma)=\int_\alpha^\beta|\gamma'(t)|\dd t.\end{gathered}
\]
The integral does not change under an increasing piecewise $C^1$ change of parameter (substitution), changes sign when the
path is reversed, and is additive under concatenation. We write $[z,w]$ for the segment from $z$ to $w$, $C(a,r)$ for the
circle $t\mapsto a+re^{it}$, $t\in[0,2\pi]$, and $\partial T$ for the boundary of a triangle $T$ traversed along its three sides.

\begin{proposition}\label{prop:pathint}
\begin{enumerate}[label=(\alph*)]
\item $\bigl|\int_\gamma f\bigr|\le L(\gamma)\max_{\gamma^*}|f|$. Hence if $f_n\to f$ uniformly on $\gamma^*$, then $\int_\gamma f_n\to\int_\gamma f$.
\item If $F$ is holomorphic on an open set containing $\gamma^*$ and $F'=f$ is continuous, then $\int_\gamma f=F(\gamma(\beta))-F(\gamma(\alpha))$. In
particular $\int_\gamma f=0$ for every closed path.
\item For $n\in\Z$, $\displaystyle\int_{C(a,r)}(z-a)^n\dd z=0$ if $n\ne-1$, and $\displaystyle\int_{C(a,r)}\frac{\mathrm dz}{z-a}=2\pi i$.
\end{enumerate}
\end{proposition}
\begin{proof}
(a) $|\int_\alpha^\beta g|\le\int_\alpha^\beta|g|$ for complex-valued $g$ (multiply by a unit complex number making the integral real). (b) On
each interval where $\gamma$ is $C^1$, $\frac{\mathrm d}{\mathrm dt}F(\gamma(t))=F'(\gamma(t))\gamma'(t)$, by the same proof as the real chain rule; integrate
and add. (c) For $n\ne-1$, $\frac{(z-a)^{n+1}}{n+1}$ is a primitive on $\C\setminus\{a\}$. For $n=-1$, $\int_0^{2\pi}\frac{ire^{it}}{re^{it}}\dd t=2\pi i$.
\end{proof}

So $\frac1z$ has no primitive on $\C\setminus\{0\}$, although it has one on the slit plane: going around the origin, the logarithm
does not come back to its starting value. Everything below grows from this observation.

\begin{theorem}[Goursat]\label{thm:goursat}
Let $f$ be holomorphic on $U$ and $T\subseteq U$ a closed triangle. Then $\int_{\partial T}f=0$.
\end{theorem}
\begin{proof}
Let $I=\int_{\partial T}f$, and let $d$ and $p$ be the diameter and perimeter of $T$. Joining the midpoints of the sides cuts $T$ into four
triangles similar to $T$ with ratio $\frac12$; oriented like $T$, their four boundary integrals add up to $I$, because each inner side
is traversed once in each direction. One of them, $T_1$, satisfies $|\int_{\partial T_1}f|\ge\frac{|I|}4$. Repeating, we get nested triangles
$T\supseteq T_1\supseteq T_2\supseteq\cdots$ with diameter $\frac d{2^k}$, perimeter $\frac p{2^k}$ and $|I|\le4^k|\int_{\partial T_k}f|$ (Figure~\ref{fig:goursat}).

The $T_k$ are nested non-empty compact sets, so they have a common point $z_0\in T\subseteq U$. As $f$ is complex differentiable at $z_0$,
$f(z)=f(z_0)+f'(z_0)(z-z_0)+\varepsilon(z)(z-z_0)$ with $\varepsilon(z)\to0$ as $z\to z_0$. The polynomial part has a primitive, so its integral over
$\partial T_k$ vanishes by Proposition~\ref{prop:pathint}(b). The rest is at most $\frac p{2^k}\cdot\frac d{2^k}\cdot\sup_{T_k}|\varepsilon|$ by (a). Hence
$|I|\le pd\sup_{T_k}|\varepsilon|$, which tends to $0$ as $k\to\infty$. So $I=0$.
\end{proof}

\begin{figure}[ht]
\centering
\fig{ch09-goursat}
\caption{Goursat's proof. Of the four midpoint triangles, keep one carrying at least a quarter of the integral, and repeat. The
triangles shrink to a point $z_0$ where $f$ is almost linear, and linear functions have integral $0$.}
\label{fig:goursat}
\end{figure}

The point of Goursat's proof (1900) is that it assumes nothing about $f'$, not even its continuity, which Cauchy had needed.
We also need a version with one exceptional point.

\begin{corollary}\label{cor:goursatpoint}
The conclusion of Theorem~\ref{thm:goursat} holds if $f$ is continuous on $U$ and holomorphic on $U\setminus\{p\}$ for one point $p$.
\end{corollary}
\begin{proof}
If $p\notin T$, this is the theorem. If $p$ is a vertex of $T$, choose points $q,q'$ on the two sides from $p$, at distance at most
$\varepsilon$ from $p$; $T$ is cut into the small triangle $pqq'$ and two triangles not containing $p$, whose integrals vanish. So $|\int_{\partial T}f|$ is
at most the perimeter of $pqq'$, which is at most $4\varepsilon$, times $\max_T|f|$; let $\varepsilon\to0$. If $p$ is any other point of $T$, cut $T$ into
two or three triangles having $p$ as a vertex.
\end{proof}

An open set $U$ is \emph{star-shaped} with respect to $c\in U$ if the segment $[c,z]$ lies in $U$ for every $z\in U$. Convex sets, such as
discs and half-planes, are star-shaped with respect to each of their points; so is the slit plane, with respect to $c=1$.

\begin{theorem}[Cauchy's theorem]\label{thm:cauchy}
Let $U$ be star-shaped with respect to $c$, and $f$ continuous on $U$ and holomorphic on $U\setminus\{p\}$. Then $F(z)=\int_{[c,z]}f$ is
holomorphic on $U$ with $F'=f$, and $\int_\gamma f=0$ for every closed path $\gamma$ in $U$.
\end{theorem}
\begin{proof}
Let $z\in U$ and $D(z,r)\subseteq U$. For $|h|<r$, every segment from $c$ to a point of $[z,z+h]$ lies in $U$, so the triangle with vertices
$c$, $z$, $z+h$ lies in $U$, and Corollary~\ref{cor:goursatpoint} gives $F(z+h)-F(z)=\int_{[z,z+h]}f$. Hence $\Delta=F(z+h)-F(z)-hf(z)$ is
\[
\Delta=\int_{[z,z+h]}(f(w)-f(z))\dd w,
\]
so by Proposition~\ref{prop:pathint}(a),
\[
\begin{aligned}{}&\Bigl|\frac{F(z+h)-F(z)}h-f(z)\Bigr|\eqbr\le\max_{w\in[z,z+h]}|f(w)-f(z)|,\end{aligned}
\]
which tends to $0$ with $h$ by continuity of $f$. Then apply Proposition~\ref{prop:pathint}(b).
\end{proof}

\begin{remark}
Cauchy's theorem holds on much more general open sets: $\int_\gamma f=0$ whenever $\gamma$ does not wind around any point outside $U$. This
\emph{homology form}, and the related Jordan curve theorem, are proved in Volume~7. The contours of this chapter are all
boundaries of convex sets, for which the star-shaped version suffices.
\end{remark}

\section{Cauchy's formula and its consequences}

\begin{theorem}[Cauchy's integral formula]\label{thm:cauchyformula}
Let $f$ be holomorphic on $U$ and $\overline D(a,r)\subseteq U$. Then for $|z-a|<r$,
\[
f(z)=\frac1{2\pi i}\int_{C(a,r)}\frac{f(w)}{w-z}\dd w.
\]
\end{theorem}
\begin{proof}
First, $\int_{C(a,r)}\frac{\mathrm dw}{w-z}=2\pi i$. Indeed, for $w$ on the circle, $\frac1{w-z}=\frac1{(w-a)-(z-a)}=\sum_{n\ge0}\frac{(z-a)^n}{(w-a)^{n+1}}$, with
terms bounded by $\frac1r\bigl(\frac{|z-a|}r\bigr)^n$; the convergence is uniform, and integrating term by term with Proposition~\ref{prop:pathint}(c)
only $n=0$ survives.

The closed disc is compact and $\C\setminus U$ is closed, so their distance is positive and $D(a,r')\subseteq U$ for some $r'>r$. Let
$g(w)=\frac{f(w)-f(z)}{w-z}$ for $w\ne z$ and $g(z)=f'(z)$: $g$ is continuous on the disc $D(a,r')$ and holomorphic except at $z$. By
Theorem~\ref{thm:cauchy}, $\int_{C(a,r)}g=0$, that is $\int_{C(a,r)}\frac{f(w)}{w-z}\dd w=f(z)\int_{C(a,r)}\frac{\mathrm dw}{w-z}=2\pi if(z)$.
\end{proof}

The values of $f$ inside the disc are averages of its values on the circle; at the centre, $f(a)=\frac1{2\pi}\int_0^{2\pi}f(a+re^{it})\dd t$
(the \emph{mean value property}). The formula has a striking consequence: the kernel $\frac1{w-z}$ is a geometric series in $z$.

\begin{theorem}[holomorphic means analytic]\label{thm:analytic}
Let $f$ be holomorphic on $U$ and $D(a,R)\subseteq U$. Then $f$ is the sum of its Taylor series on $D(a,R)$:
\[
\begin{gathered}f(z)=\sum_{n\ge0}c_n(z-a)^n,\qq c_n=\frac{f^{(n)}(a)}{n!}=\frac1{2\pi i}\int_{C(a,r)}\frac{f(w)\dd w}{(w-a)^{n+1}}\end{gathered}
\]
for every $r\in\left]0,R\right[$.
In particular $f$ has derivatives of all orders, all holomorphic on $U$.
\end{theorem}
\begin{proof}
Fix $0<r<R$ and $|z-a|<r$. Expanding $\frac1{w-z}$ as in the previous proof, uniformly for $w\in C(a,r)$, and integrating term by term
in Cauchy's formula gives $f(z)=\sum c_n(r)(z-a)^n$, with $c_n(r)$ the integral in the statement. So $f$ is a power series on $D(a,r)$, and
Theorem~\ref{thm:powerseries} gives $c_n(r)=\frac{f^{(n)}(a)}{n!}$, independent of $r$. Every $z\in D(a,R)$ is in some $D(a,r)$ with $r<R$. Finally, $f'$ is
holomorphic near each point of $U$, being there the sum of a power series.
\end{proof}

\begin{corollary}[Cauchy's estimates]\label{cor:radius}
\begin{enumerate}[label=(\alph*)]
\item With $M_r=\max_{C(a,r)}|f|$, $\ \bigl|f^{(n)}(a)\bigr|\le\frac{n!\,M_r}{r^n}$.
\item The Taylor series of $f$ at $a$ has radius of convergence at least the radius of the largest open disc around $a$ contained in
$U$. If $f$ is unbounded near a point $b$ of the boundary of that disc, the radius is exactly $|b-a|$.
\end{enumerate}
\end{corollary}
\begin{proof}
(a) Bound the integral for $c_n$ by Proposition~\ref{prop:pathint}(a): $|c_n|\le\frac1{2\pi}\cdot2\pi r\cdot\frac{M_r}{r^{n+1}}$. (b) The first part is
Theorem~\ref{thm:analytic}. If the radius were larger than $|b-a|$, the sum of the series would be continuous, hence bounded, near $b$,
and it equals $f$ on $D(a,|b-a|)$; so $f$ would be bounded near $b$ on that disc.
\end{proof}

\begin{example}
This answers the think box: $\frac1{1+z^2}$ is holomorphic except at $\pm i$, where it is unbounded, so its Taylor series at $0$
has radius $1$, the distance to $\pm i$. The real line simply cannot see the obstruction. Similarly $\frac z{e^z-1}=\frac1{h(z)}$ with
$h(z)=\sum_{n\ge0}\frac{z^n}{(n+1)!}$ entire; $h$ vanishes exactly where $e^z=1$ and $z\ne0$, that is on $2\pi i\Z\setminus\{0\}$. So $\frac z{e^z-1}$ is holomorphic on $D(0,2\pi)$
and unbounded near $\pm2\pi i$: the generating function of the Bernoulli numbers has radius exactly $2\pi$, as Chapters~1 and~3
announced.
\end{example}

A continuous function with vanishing triangle integrals is holomorphic; this converse of Goursat makes holomorphy
stable under limits and integrals.

\begin{theorem}[Morera]\label{thm:morera}
Let $f$ be continuous on $U$ with $\int_{\partial T}f=0$ for every closed triangle $T\subseteq U$. Then $f$ is holomorphic on $U$.
\end{theorem}
\begin{proof}
On a disc $D\subseteq U$ centred at $c$, the proof of Theorem~\ref{thm:cauchy} (which used only the vanishing of triangle integrals)
shows that $F(z)=\int_{[c,z]}f$ is holomorphic with $F'=f$. By Theorem~\ref{thm:analytic}, $F'$ is holomorphic on $D$.
\end{proof}

\begin{theorem}[limits and integrals]\label{thm:holoparam}
\begin{enumerate}[label=(\alph*)]
\item (Weierstrass) If $f_n$ are holomorphic on $U$ and converge to $f$ uniformly on every closed disc contained in $U$, then $f$ is
holomorphic, and $f_n'\to f'$ uniformly on every closed disc contained in $U$. In particular a series of holomorphic functions converging
uniformly on such discs may be differentiated term by term.
\item Let $I$ be an interval and $f$ continuous on $I\times U$, holomorphic in $z$ for each $t\in I$. Suppose that for every closed disc
$K\subseteq U$ there is a continuous integrable $g_K$ on $I$ with $|f(t,z)|\le g_K(t)$ for $z\in K$. Then $F(z)=\int_If(t,z)\dd t$ is holomorphic on $U$.
\end{enumerate}
\end{theorem}
\begin{proof}
(a) $f$ is continuous, as a locally uniform limit of continuous functions. For a triangle $T$ in a disc $D\subseteq U$, $\int_{\partial T}f=\lim\int_{\partial T}f_n=0$ by
Proposition~\ref{prop:pathint}(a) and Goursat, so $f$ is holomorphic by Morera. If $\overline D(a,r)\subseteq U$, choose $\rho>0$ with $\overline D(a,r+2\rho)\subseteq U$. For
$|z-a|\le r$, Theorem~\ref{thm:analytic} on the disc $D(z,\rho)$ gives $f_n'(z)-f'(z)=\frac1{2\pi i}\int_{C(z,\rho)}\frac{f_n(w)-f(w)}{(w-z)^2}\dd w$, whose modulus is at most
$\frac1\rho\max_{\overline D(a,r+\rho)}|f_n-f|$, which tends to $0$ uniformly in $z$.

(b) Continuity: if $z_n\to z$ inside a closed disc $K\subseteq U$, then $f(t,z_n)\to f(t,z)$ uniformly for $t$ in every closed bounded subinterval
(uniform continuity on compact sets), with domination by $g_K$, so $F(z_n)\to F(z)$ by Corollary~\ref{cor:dct}. Next let $T$ be a triangle in a disc
contained in $U$, with $\partial T$ made of three segments parametrised on $[0,1]$ and $K$ a closed disc containing $T$. For a closed bounded
$[u,v]\subseteq I$, Fubini on the rectangle $[u,v]\times[0,1]$ (Theorem~\ref{thm:fubinirect}) gives $\int_{\partial T}\int_u^vf(t,z)\dd t\dd z=\int_u^v\int_{\partial T}f(t,z)\dd z\dd t=0$ by
Goursat. As $[u,v]$ grows to $I$, the left side tends to $\int_{\partial T}F$, the difference being at most $L(\partial T)\int_{I\setminus[u,v]}g_K$. So $\int_{\partial T}F=0$,
and Morera applies.
\end{proof}

\subsection*{Liouville, zeros and the maximum principle}

\begin{theorem}[Liouville]\label{thm:liouville}
A bounded entire function is constant.
\end{theorem}
\begin{proof}
If $|f|\le M$ on $\C$, Theorem~\ref{thm:analytic} on $D(0,R)$ for every $R$ gives $f(z)=\sum c_nz^n$ on $\C$, and Corollary~\ref{cor:radius}(a) gives $|c_n|\le\frac M{r^n}$ for
every $r>0$. Letting $r\to\infty$, $c_n=0$ for $n\ge1$.
\end{proof}

\begin{corollary}[fundamental theorem of algebra]\label{cor:fta}
Every non-constant complex polynomial has a root in $\C$.
\end{corollary}
\begin{proof}
Let $p(z)=a_nz^n+\dots+a_0$ with $n\ge1$ and $a_n\ne0$. Then $|p(z)|\ge|a_n||z|^n\bigl(1-\sum_{k<n}\frac{|a_k|}{|a_n||z|^{n-k}}\bigr)\to\infty$ as $|z|\to\infty$. If $p$
had no root, $\frac1p$ would be entire, tend to $0$ at infinity, hence be bounded (it is continuous on a large closed disc and small outside),
and so constant by Liouville: a contradiction.
\end{proof}

Volume~0 (Theorem~3.3) proved this by d'Alembert's argument on the minimum of $|p|$. Here it falls out of Liouville in two
lines, and Section~\ref{sec:argument} gives a third proof that also counts the roots.

An open set $U$ is \emph{connected} if any two of its points are joined by a path in $U$.

\begin{theorem}[zeros and the identity theorem]\label{thm:identity}
Let $f$ be holomorphic on a connected open set $U$, and not identically zero.
\begin{enumerate}[label=(\alph*)]
\item For every $a\in U$ there are $m\ge0$ and $g$ holomorphic on $U$ with $g(a)\ne0$ and $f(z)=(z-a)^mg(z)$. In particular the zeros of $f$ are isolated.
\item Consequently, two holomorphic functions on $U$ that agree on a set having an accumulation point in $U$ are equal.
\end{enumerate}
\end{theorem}
\begin{proof}
Let $A$ be the set of $z\in U$ where all derivatives $f^{(n)}(z)$, $n\ge0$, vanish. $A$ is open: near $z\in A$, $f$ is the sum of its Taylor series
(Theorem~\ref{thm:analytic}), which is $0$, so $f$ and its derivatives vanish on a disc. $U\setminus A$ is open: if $f^{(n)}(z)\ne0$, then $f^{(n)}\ne0$ near $z$ by
continuity. If both $A$ and $U\setminus A$ were non-empty, a path $\gamma:[0,1]\to U$ from a point of $A$ to a point of $U\setminus A$ would give a partition of $[0,1]$
into the two non-empty disjoint open sets $\gamma^{-1}(A)$ and $\gamma^{-1}(U\setminus A)$, which is impossible since $[0,1]$ is connected (take the supremum of the
first set that contains $0$). As $f\ne0$, $A$ is empty.

(a) At $a$, let $m$ be the least $n$ with $f^{(n)}(a)\ne0$. Near $a$, $f(z)=\sum_{n\ge m}c_n(z-a)^n=(z-a)^m\sum_{k\ge0}c_{m+k}(z-a)^k$; set
$g(z)=\frac{f(z)}{(z-a)^m}$ for $z\ne a$ and $g(a)=c_m\ne0$. Then $g$ is holomorphic (near $a$ it is a power series), and $g\ne0$ near $a$ by continuity,
so $a$ is the only zero of $f$ in a small disc. (b) Apply (a) to the difference: an accumulation point of zeros is not an isolated zero.
\end{proof}

So identities between holomorphic functions proved on a real interval hold wherever both sides are holomorphic on a
connected open set: $\sin^2z+\cos^2z=1$ on $\C$, for instance. Section~\ref{sec:gammac} uses this for $\Gamma$.

\begin{theorem}[maximum principle]\label{thm:maxprinciple}
If $f$ is holomorphic on a connected open set $U$ and $|f|$ has a local maximum at $a\in U$, then $f$ is constant.
\end{theorem}
\begin{proof}
Let $|f|\le|f(a)|$ on $D(a,R)$, and write $f(a)=|f(a)|e^{i\alpha}$. For $0<r<R$, the mean value property gives
$|f(a)|=\frac1{2\pi}\int_0^{2\pi}\operatorname{Re}\bigl(e^{-i\alpha}f(a+re^{it})\bigr)\dd t$, while the integrand is at most $|f(a+re^{it})|\le|f(a)|$. A continuous function
bounded by a constant with average equal to that constant is the constant, so $\operatorname{Re}(e^{-i\alpha}f)=|f(a)|\ge|f|$ on the circle, which forces
$e^{-i\alpha}f=|f(a)|$ there. So $f=f(a)$ on $D(a,R)$, and on $U$ by Theorem~\ref{thm:identity}.
\end{proof}

On a bounded open set, $|f|$ continuous on the closure therefore reaches its maximum on the boundary: the same averaging
argument as the ``maximum principle in miniature'' of Chapter~3.

\section{Singularities, Laurent series and residues}

Let $f$ be holomorphic on a punctured disc $D(a,r)\setminus\{a\}$: $a$ is an \emph{isolated singularity}.

\begin{theorem}[Riemann's removable singularity theorem]\label{thm:removable}
If $f$ is bounded near $a$, it extends to a holomorphic function on $D(a,r)$.
\end{theorem}
\begin{proof}
Let $h(z)=(z-a)^2f(z)$ for $z\ne a$ and $h(a)=0$. Then $\frac{h(z)-h(a)}{z-a}=(z-a)f(z)\to0$, so $h$ is holomorphic on $D(a,r)$ with $h(a)=h'(a)=0$. By
Theorem~\ref{thm:analytic}, $h(z)=\sum_{n\ge2}c_n(z-a)^n$, and $\sum_{n\ge2}c_n(z-a)^{n-2}$ is the required extension of $f$.
\end{proof}

\begin{lemma}\label{lem:annulus}
Let $g$ be holomorphic on the annulus $r<|z-a|<R$. Then $\int_{C(a,\rho)}g$ does not depend on $\rho\in\left]r,R\right[$.
\end{lemma}
\begin{proof}
Let $\phi(\rho,t)=g(a+\rho e^{it})\,i\rho e^{it}$, so the integral is $J(\rho)=\int_0^{2\pi}\phi(\rho,t)\dd t$. Since $g'$ is continuous (Theorem~\ref{thm:analytic}), the chain
rule gives
\[
\begin{aligned}\partial_\rho\phi&=g'(a+\rho e^{it})\,i\rho e^{2it}+g(a+\rho e^{it})\,ie^{it}\eqbr=\partial_t\bigl(g(a+\rho e^{it})e^{it}\bigr).\end{aligned}
\]
By Theorem~\ref{thm:paramcompact}, $J'(\rho)=\int_0^{2\pi}\partial_t\bigl(g(a+\rho e^{it})e^{it}\bigr)\dd t=0$, the function in brackets being $2\pi$-periodic in $t$.
\end{proof}

\begin{theorem}[Laurent]\label{thm:laurent}
Let $f$ be holomorphic on the annulus $A=\{r<|z-a|<R\}$, with $0\le r<R\le\infty$. Then
\[
\begin{gathered}f(z)=\sum_{n\in\Z}c_n(z-a)^n,\qq c_n=\frac1{2\pi i}\int_{C(a,\rho)}\frac{f(w)\dd w}{(w-a)^{n+1}}\quad(r<\rho<R),\end{gathered}
\]
where the coefficients do not depend on $\rho$. The part $\sum_{n\ge0}$ converges on $D(a,R)$; the part $\sum_{n<0}$ converges for $|z-a|>r$ and is
holomorphic there; both converge uniformly on every closed annulus $\rho_1\le|z-a|\le\rho_2$ inside $A$. The coefficients are unique.
\end{theorem}
\begin{proof}
Independence of $\rho$ is Lemma~\ref{lem:annulus} applied to $\frac{f(w)}{(w-a)^{n+1}}$. Fix $z$ and $r<\rho_1<|z-a|<\rho_2<R$. The function
$g(w)=\frac{f(w)-f(z)}{w-z}$ is bounded near $z$, hence holomorphic on $A$ by Theorem~\ref{thm:removable}, and Lemma~\ref{lem:annulus} gives
$\int_{C(a,\rho_2)}g=\int_{C(a,\rho_1)}g$. Now $\int_{C(a,\rho_2)}\frac{\mathrm dw}{w-z}=2\pi i$ (proof of Theorem~\ref{thm:cauchyformula}), while $\int_{C(a,\rho_1)}\frac{\mathrm dw}{w-z}=0$
because there $\frac1{w-z}=-\sum_{n\ge0}\frac{(w-a)^n}{(z-a)^{n+1}}$, uniformly, and each term has a primitive. Therefore
\[
\begin{aligned}f(z)&=\frac1{2\pi i}\int_{C(a,\rho_2)}\frac{f(w)\dd w}{w-z}\eqbr\quad-\frac1{2\pi i}\int_{C(a,\rho_1)}\frac{f(w)\dd w}{w-z}.\end{aligned}
\]
Expanding the kernels in the two geometric series just written and integrating term by term gives $\sum_{n\ge0}c_n(z-a)^n$ for the first
integral and $\sum_{m\ge1}c_{-m}(z-a)^{-m}$ for the second. In the variable $u=\frac1{z-a}$, the second part is a power series converging for $|u|<\frac1{\rho_1}$ for
every $\rho_1>r$, so Lemma~\ref{lem:abel} and Theorem~\ref{thm:powerseries} give its convergence and holomorphy; the first part is handled likewise.
For uniqueness, if $f=\sum d_n(z-a)^n$ converges uniformly on $C(a,\rho)$, multiplying by $(z-a)^{-k-1}$ and integrating term by term
gives $2\pi i\,d_k$ by Proposition~\ref{prop:pathint}(c), so $d_k=c_k$.
\end{proof}

At an isolated singularity ($r=0$), the \emph{principal part} $P(z)=\sum_{n<0}c_n(z-a)^n$ is holomorphic on $\C\setminus\{a\}$, and $f-P$ extends
holomorphically to $a$. The singularity is \emph{removable} if $P=0$, a \emph{pole of order $m$} if $P$ ends at $c_{-m}\ne0$, and \emph{essential}
if $P$ has infinitely many terms, as for $e^{1/z}=\sum_{n\ge0}\frac{z^{-n}}{n!}$ at $0$.

\begin{definition}
The \emph{residue} of $f$ at the isolated singularity $a$ is $\operatorname{Res}_af=c_{-1}=\frac1{2\pi i}\int_{C(a,\rho)}f$ for small $\rho$.
\end{definition}

\begin{proposition}[computing residues]\label{prop:rescalc}
\begin{enumerate}[label=(\alph*)]
\item If $a$ is a pole of order at most $m$, then $\operatorname{Res}_af=\frac1{(m-1)!}\lim_{z\to a}\frac{\mathrm d^{m-1}}{\mathrm dz^{m-1}}\bigl((z-a)^mf(z)\bigr)$; for a simple pole,
$\operatorname{Res}_af=\lim_{z\to a}(z-a)f(z)$.
\item If $g$, $h$ are holomorphic near $a$ with $h(a)=0\ne h'(a)$, then $\operatorname{Res}_a\frac gh=\frac{g(a)}{h'(a)}$.
\end{enumerate}
\end{proposition}
\begin{proof}
(a) $(z-a)^mf(z)=\sum_{n\ge-m}c_n(z-a)^{n+m}$ is a power series, whose coefficient of $(z-a)^{m-1}$ is $c_{-1}$; apply Theorem~\ref{thm:powerseries}. (b) The pole is
at most simple, and $(z-a)\frac{g(z)}{h(z)}=g(z)\frac{z-a}{h(z)-h(a)}\to\frac{g(a)}{h'(a)}$.
\end{proof}

For instance $\pi\cot\pi z=\frac{\pi\cos\pi z}{\sin\pi z}$ has at each integer $n$ a simple pole with residue $\frac{\pi\cos\pi n}{\pi\cos\pi n}=1$. This is the reason it
appears in Chapter~3 and again below: it is a machine with residue $1$ at every integer.

\section{The residue theorem}

For a closed path $\gamma$ and $a\notin\gamma^*$, the \emph{winding number} is $n(\gamma,a)=\frac1{2\pi i}\int_\gamma\frac{\mathrm dz}{z-a}$.

\begin{proposition}\label{prop:winding}
$n(\gamma,a)$ is an integer: the number of times $\gamma$ turns around $a$.
\end{proposition}
\begin{proof}
Let $\gamma:[\alpha,\beta]\to\C$ and $\phi(t)=\int_\alpha^t\frac{\gamma'(s)}{\gamma(s)-a}\dd s$. Then $\frac{\mathrm d}{\mathrm dt}\bigl((\gamma(t)-a)e^{-\phi(t)}\bigr)=0$ wherever $\gamma$ is $C^1$, so this
continuous function is constant, and $\gamma(\beta)=\gamma(\alpha)$ gives $e^{\phi(\beta)}=1$, so $\phi(\beta)\in2\pi i\Z$. As $e^{\phi(t)}=\frac{\gamma(t)-a}{\gamma(\alpha)-a}$, the imaginary part
of $\phi$ is a continuous choice of the angle of $\gamma(t)-a$, and $\frac{\operatorname{Im}\phi(\beta)}{2\pi}$ counts its full turns.
\end{proof}

Our contours are boundaries of bounded convex open sets $D$: discs, half-discs, rectangles, squares, circular sectors of angle
at most $\pi$. For each point $b\in D$, every ray from $b$ meets the boundary exactly once, at $b+\rho(\theta)e^{i\theta}$, where $\rho$ is given on each
side or arc by an elementary piecewise $C^1$ formula. We say that a closed path $\gamma$ is the \emph{positive boundary} of $D$ if, for
every $b\in D$, it is a change of parameter of $\theta\mapsto b+\rho(\theta)e^{i\theta}$, $\theta\in[\theta_0,\theta_0+2\pi]$; this is the usual counterclockwise contour. Then
\[
\begin{aligned}2\pi i\,n(\gamma,b)&=\int_{\theta_0}^{\theta_0+2\pi}\frac{\rho'(\theta)e^{i\theta}+i\rho(\theta)e^{i\theta}}{\rho(\theta)e^{i\theta}}\dd\theta\eqbr=\bigl[\log\rho\bigr]_{\theta_0}^{\theta_0+2\pi}+2\pi i=2\pi i,\end{aligned}
\]
so $n(\gamma,b)=1$ for every $b\in D$.

\begin{theorem}[residue theorem]\label{thm:residue}
Let $D$ be a bounded convex open set with positive boundary $\gamma$, and $U$ an open set containing $\overline D$. Let $f$ be holomorphic on
$U\setminus\{a_1,\dots,a_k\}$, where $a_1,\dots,a_k\in D$. Then
\[
\int_\gamma f(z)\dd z=2\pi i\sum_{j=1}^k\operatorname{Res}_{a_j}f.
\]
\end{theorem}
\begin{proof}
Let $P_j$ be the principal part of $f$ at $a_j$, holomorphic on $\C\setminus\{a_j\}$ (Theorem~\ref{thm:laurent}). Near $a_j$, $f-\sum_iP_i$ is $(f-P_j)$, which extends
holomorphically, minus the $P_i$ with $i\ne j$, which are holomorphic at $a_j$. So $g=f-\sum_iP_i$ extends to a holomorphic function on $U$.
For small $\varepsilon>0$, the set $D_\varepsilon$ of points at distance less than $\varepsilon$ from $\overline D$ is open, convex (if $|z-d|<\varepsilon$ and $|z'-d'|<\varepsilon$ with
$d,d'\in\overline D$, the same holds for convex combinations) and contained in $U$ (the compact $\overline D$ is at positive distance from $\C\setminus U$).
By Cauchy's theorem on $D_\varepsilon$, $\int_\gamma g=0$.

For each $j$, $\gamma^*$ stays at distance at least some $\delta>0$ from $a_j$, where the series $P_j(z)=\sum_{m\ge1}c_{-m}(z-a_j)^{-m}$ converges uniformly
(Theorem~\ref{thm:laurent}). Integrating term by term, the terms with $m\ge2$ have primitives and give $0$, and the term $m=1$ gives
$c_{-1}\cdot2\pi i\,n(\gamma,a_j)=2\pi i\operatorname{Res}_{a_j}f$.
\end{proof}

\begin{figure}[ht]
\centering
\fig{ch09-contours}
\caption{The three contours of this chapter, with the singularities they enclose. Left: a half-disc of radius $R$, for integrals of
rational functions over $\R$. Middle: a rectangle of height $2\pi$, for integrands periodic under $t\mapsto t+2\pi i$ up to a factor. Right:
the square $C_N$ with vertices $(N+\frac12)(\pm1\pm i)$, which avoids the integers, the poles of $\pi\cot\pi z$.}
\label{fig:contours}
\end{figure}

\section{Real integrals by residues}

The method: extend the integrand to the complex plane, close the real segment $[-R,R]$ by a path on which the integral is
small or is a known multiple of the one wanted, and let $R\to\infty$.

\begin{proposition}\label{prop:rationalint}
Let $P$, $Q$ be polynomials with $\deg Q\ge\deg P+2$ and no real zero of $Q$. Then
\[
\int_{-\infty}^\infty\frac{P(x)}{Q(x)}\dd x=2\pi i\sum_{\operatorname{Im}a>0}\operatorname{Res}_a\frac PQ,
\]
the sum running over the zeros of $Q$ in the upper half-plane.
\end{proposition}
\begin{proof}
There are $C$ and $R_0$ with $\bigl|\frac{P(z)}{Q(z)}\bigr|\le\frac C{|z|^2}$ for $|z|\ge R_0$, so the integral converges absolutely. For $R>R_0$ larger than the moduli of
all zeros of $Q$, apply Theorem~\ref{thm:residue} to the half-disc of radius $R$ (Figure~\ref{fig:contours}, left): the integral over $[-R,R]$ plus the
integral over the arc $\Gamma_R$ equals $2\pi i$ times the sum of residues. By Proposition~\ref{prop:pathint}(a), $\bigl|\int_{\Gamma_R}\frac PQ\bigr|\le\pi R\cdot\frac C{R^2}\to0$.
\end{proof}

\begin{example}
$\int_{-\infty}^\infty\frac{\mathrm dx}{1+x^4}=\frac\pi{\sqrt2}$. The poles in the upper half-plane are $a=e^{i\pi/4}$ and $e^{3i\pi/4}$. By Proposition~\ref{prop:rescalc}(b) the
residue at a root $a$ of $z^4=-1$ is $\frac1{4a^3}=\frac a{4a^4}=-\frac a4$, and $e^{i\pi/4}+e^{3i\pi/4}=i\sqrt2$. So the integral is $2\pi i\cdot\bigl(-\frac{i\sqrt2}4\bigr)=\frac{\pi\sqrt2}2$.
\end{example}

\begin{example}[the transform of $\frac1{1+x^2}$, promised in Chapter~8]
For $\xi\le0$ and $\operatorname{Im}z\ge0$, $|e^{-2\pi iz\xi}|=e^{2\pi\xi\operatorname{Im}z}\le1$. So the proof of Proposition~\ref{prop:rationalint} works for $f(z)=\frac{e^{-2\pi iz\xi}}{1+z^2}$, the arc
still contributing at most $\frac{\pi R}{R^2-1}$. The only pole in the upper half-plane is $i$, with residue $\frac{e^{2\pi\xi}}{2i}$, so the transform is $\pi e^{2\pi\xi}=\pi e^{-2\pi|\xi|}$.
For $\xi\ge0$ the function is even in $\xi$ (substitute $x\mapsto-x$). This is Exercise~\ref{ex:lorentz}, without the inversion theorem.
\end{example}

\begin{example}[the reflection formula again]\label{ex:rectangle}
For $0<a<1$, $\displaystyle\int_{-\infty}^\infty\frac{e^{at}}{1+e^t}\dd t=\frac\pi{\sin\pi a}$. The integrand $f$ satisfies $f(t+2\pi i)=e^{2\pi ia}f(t)$, which suggests the rectangle
with vertices $\pm R$ and $\pm R+2\pi i$ (Figure~\ref{fig:contours}, middle). Inside, $1+e^t$ vanishes only at $i\pi$, with residue $\frac{e^{i\pi a}}{e^{i\pi}}=-e^{i\pi a}$ by
Proposition~\ref{prop:rescalc}(b). On the right side $|f|\le\frac{e^{aR}}{e^R-1}$, on the left side $|f|\le\frac{e^{-aR}}{1-e^{-R}}$; both sides have length $2\pi$ and their
integrals tend to $0$ because $0<a<1$. The top side, traversed from right to left, gives $-e^{2\pi ia}$ times the integral $I_R$ over $[-R,R]$. So
\[
\begin{aligned}(1-e^{2\pi ia})I&=-2\pi ie^{i\pi a},\eqbr I=\frac{2\pi i}{e^{i\pi a}-e^{-i\pi a}}=\frac\pi{\sin\pi a}.\end{aligned}
\]
The substitution $x=e^t$ turns $I$ into $\int_0^\infty\frac{x^{a-1}}{1+x}\dd x$, which is $B(a,1-a)=\Gamma(a)\Gamma(1-a)$ (substitute $x=\frac u{1-u}$ in the Beta integral of
Theorem~\ref{thm:betagamma}): a second proof of Theorem~\ref{thm:reflection}, this time without the sine product.
\end{example}

\begin{example}[the Gaussian, a third time]
For real $\xi$, $f(z)=e^{-\pi z^2}$ is entire, so its integral over the rectangle with vertices $\pm R$, $\pm R+i\xi$ is $0$. On the vertical sides,
$|e^{-\pi(\pm R+iy)^2}|=e^{-\pi R^2+\pi y^2}\le e^{-\pi R^2+\pi\xi^2}$, so their integrals tend to $0$. Hence $\int_{\R}e^{-\pi(x+i\xi)^2}\dd x=\int_{\R}e^{-\pi x^2}\dd x=1$, and
expanding $e^{-\pi(x+i\xi)^2}=e^{-\pi x^2}e^{-2\pi ix\xi}e^{\pi\xi^2}$ gives $\hat g(\xi)=e^{-\pi\xi^2}$: Theorem~\ref{thm:gaussianft} by shifting the line of integration.
\end{example}

\begin{beauty}
In these three examples the real integral never sees $i$. Yet the answer is decided by a single point of the complex plane
($e^{i\pi/4}$, then $i$, then $i\pi$), and the rest of the integrand is irrelevant. This is the answer to ``how can a real
integral know about $i$?'': the function is rigid, so its values on the line determine it everywhere, and the residue is
where that information is concentrated.
\end{beauty}

\section{Sums by residues}

The function $\pi\cot\pi z$ has residue $1$ at every integer, so $\pi\cot(\pi z)f(z)$ has residue $f(n)$ at every integer $n$ where $f$ is
holomorphic. Integrating over large squares turns a sum over the integers into a few residues of $f$.

\begin{lemma}\label{lem:cotbound}
On the square $C_N$ with vertices $(N+\frac12)(\pm1\pm i)$, $N\ge0$, we have $|\cot\pi z|\le\coth\frac\pi2<2$.
\end{lemma}
\begin{proof}
$\cot\pi z=i\frac{e^{2\pi iz}+1}{e^{2\pi iz}-1}$. On the horizontal sides, $|\operatorname{Im}z|=N+\frac12$, so $q=|e^{2\pi iz}|=e^{-2\pi\operatorname{Im}z}$ is at most $e^{-\pi}$ or at least $e^\pi$; in both cases
$\bigl|\frac{q'+1}{q'-1}\bigr|\le\frac{q+1}{|q-1|}\le\frac{1+e^{-\pi}}{1-e^{-\pi}}=\coth\frac\pi2$ for $|q'|=q$. On the vertical sides, $z=\pm(N+\frac12)+iy$ and $\cot\pi z=-\tan(i\pi y)=-i\tanh\pi y$,
of modulus less than $1$. Finally $\coth\frac\pi2\approx1.09$.
\end{proof}

\begin{theorem}\label{thm:residuesum}
Let $f$ be holomorphic on $\C$ except at a finite set $S$, with $|f(z)|\le\frac C{|z|^2}$ for $|z|\ge R_0$. Then
\[
\sum_{n\in\Z\setminus S}f(n)=-\sum_{a\in S}\operatorname{Res}_a\bigl(\pi\cot(\pi z)f(z)\bigr),
\]
the series converging absolutely.
\end{theorem}
\begin{proof}
Absolute convergence follows from the bound. Let $N$ be large enough that $N+\frac12>R_0$ and $S$ lies inside $C_N$. The function
$\pi\cot(\pi z)f(z)$ is holomorphic on a neighbourhood of the closed square except at the integers $n$ with $|n|\le N$ and at the points of
$S$; at an integer $n\notin S$ it has a simple pole with residue $f(n)$ (Proposition~\ref{prop:rescalc}(a)). By Theorem~\ref{thm:residue}, with
$g(z)=\pi\cot(\pi z)f(z)$,
\[
\frac1{2\pi i}\int_{C_N}g(z)\dd z=\sum_{|n|\le N,\ n\notin S}f(n)+\sum_{a\in S}\operatorname{Res}_ag.
\]
By Lemma~\ref{lem:cotbound}, the left side is at most $\frac1{2\pi}\cdot8\bigl(N+\frac12\bigr)\cdot2\pi\cdot\frac C{(N+1/2)^2}$, which tends to $0$.
\end{proof}

We need the start of the Laurent series of $\pi\cot\pi z$ at $0$. From $\sin w=w\bigl(1-\frac{w^2}6+O(w^4)\bigr)$ and $\cos w=1-\frac{w^2}2+O(w^4)$,
$w\cot w=\bigl(1-\frac{w^2}2+\cdots\bigr)\bigl(1+\frac{w^2}6+\cdots\bigr)=1-\frac{w^2}3+O(w^4)$, an even power series; so
\[
\pi\cot\pi z=\frac1z-\frac{\pi^2}3z+O(z^3).
\]

\begin{corollary}[Basel, the sixth proof]\label{cor:baselres}
$\displaystyle\sum_{n\ge1}\frac1{n^2}=\frac{\pi^2}6$.
\end{corollary}
\begin{proof}
Take $f(z)=\frac1{z^2}$ and $S=\{0\}$. The Laurent series $\frac{\pi\cot\pi z}{z^2}=\frac1{z^3}-\frac{\pi^2}3\cdot\frac1z+O(z)$ has residue $-\frac{\pi^2}3$, so $\sum_{n\ne0}\frac1{n^2}=\frac{\pi^2}3$.
\end{proof}

\begin{corollary}[Euler's partial fractions, on $\C$]\label{cor:cotc}
For every $z\in\C\setminus\Z$, $\displaystyle\pi\cot\pi z=\frac1z+\sum_{n\ge1}\frac{2z}{z^2-n^2}$.
\end{corollary}
\begin{proof}
Fix $z\notin\Z$ and apply Theorem~\ref{thm:residuesum} to $f(w)=\frac1{w(w-z)}$, in the variable $w$, with $S=\{0,z\}$. At $w=z$ the residue of
$\frac{\pi\cot\pi w}{w(w-z)}$ is $\frac{\pi\cot\pi z}z$. At $w=0$, multiply $\pi\cot\pi w=\frac1w+O(w)$ by $\frac1w\cdot\frac1{w-z}=\frac1w\bigl(-\frac1z-\frac w{z^2}+O(w^2)\bigr)$: the coefficient of $\frac1w$ is
$-\frac1{z^2}$. So $\sum_{n\ne0}\frac1{n(n-z)}=\frac1{z^2}-\frac{\pi\cot\pi z}z$. Multiply by $-z$ and use $\frac{-z}{n(n-z)}=\frac1{z-n}+\frac1n$; grouping $n$ with $-n$,
the terms $\frac1n$ cancel and $\frac1{z-n}+\frac1{z+n}=\frac{2z}{z^2-n^2}$.
\end{proof}

Chapter~3 proved this for real $x$ with Herglotz's trick; now both sides are holomorphic on $\C\setminus\Z$, and residues give
the identity there directly. The same method with $f(z)=z^{-2k}$ gives $\zeta(2k)=-\frac12\operatorname{Res}_0\frac{\pi\cot\pi z}{z^{2k}}$, half the coefficient of $z^{2k-1}$ in
$-\pi\cot\pi z$; with the Bernoulli expansion $\pi z\cot\pi z=\sum_{k\ge0}\frac{(-1)^kB_{2k}(2\pi z)^{2k}}{(2k)!}$ (from $\frac t{e^t-1}+\frac t2=\frac t2\coth\frac t2$ at $t=2\pi iz$), this is
Euler's formula of Chapter~3 once more.

\section{The Gamma function in the complex plane}\label{sec:gammac}

\begin{theorem}\label{thm:gammac}
\begin{enumerate}[label=(\alph*)]
\item $\Gamma(z)=\int_0^\infty t^{z-1}e^{-t}\dd t$ is holomorphic on $\operatorname{Re}z>0$, with $|\Gamma(z)|\le\Gamma(\operatorname{Re}z)$, and $\Gamma(z+1)=z\Gamma(z)$.
\item $\Gamma$ extends to a holomorphic function on the complement of $\{0,-1,-2,\dots\}$, with simple poles at $-n$ of residue $\frac{(-1)^n}{n!}$, satisfying
$\Gamma(z+1)=z\Gamma(z)$ everywhere.
\item $\Gamma(z)\Gamma(1-z)=\frac\pi{\sin\pi z}$ for $z\in\C\setminus\Z$. Hence $\Gamma$ has no zeros, and $\frac1\Gamma$ extends to an entire function vanishing exactly at
$0,-1,-2,\dots$
\item For real $y$, $|\Gamma(\frac12+iy)|^2=\frac\pi{\cosh\pi y}$.
\end{enumerate}
\end{theorem}
\begin{proof}
(a) Here $t^{z-1}=e^{(z-1)\log t}$, so $|t^{z-1}e^{-t}|=t^{x-1}e^{-t}$ with $x=\operatorname{Re}z$; this gives the bound, and on a closed disc in $x_1\le\operatorname{Re}z\le x_2$ (with
$x_1>0$) the integrand is dominated by $t^{x_1-1}e^{-t}+t^{x_2-1}e^{-t}$, integrable on $\left]0,\infty\right[$. Theorem~\ref{thm:holoparam}(b) gives holomorphy. Both sides of
$\Gamma(z+1)=z\Gamma(z)$ are holomorphic on the connected half-plane and agree on $\left]0,\infty\right[$ (Proposition~\ref{prop:gammafe}), so they agree everywhere by
Theorem~\ref{thm:identity}.

(b) For $n\ge0$, $\Gamma_n(z)=\frac{\Gamma(z+n+1)}{z(z+1)\cdots(z+n)}$ is holomorphic on $\operatorname{Re}z>-n-1$ except at $0,-1,\dots,-n$, and equals $\Gamma$ on $\operatorname{Re}z>0$ by
the functional equation; $\Gamma_{n+1}=\Gamma_n$ where both are defined, by the same identity. Together they define the extension, which
satisfies the functional equation by the identity theorem (on the connected set $\C\setminus\{0,-1,\dots\}$). Near $-n$, $(z+n)\Gamma(z)=\frac{\Gamma(z+n+1)}{z(z+1)\cdots(z+n-1)}$
tends to $\frac{\Gamma(1)}{(-n)(-n+1)\cdots(-1)}=\frac{(-1)^n}{n!}$, so the pole is simple with that residue.

(c) Both sides are holomorphic on $\C\setminus\Z$, which is connected, and agree on $\left]0,1\right[$ by Theorem~\ref{thm:reflection}. If $\Gamma(z_0)=0$, then $z_0\notin\Z$
(at the positive integers $\Gamma=(n-1)!$), and the left side would vanish at $z_0$ while the right side never does. So $\frac1\Gamma$ is holomorphic
except at the poles of $\Gamma$, where it tends to $0$; by Theorem~\ref{thm:removable}, it is entire.

(d) $\overline{\Gamma(z)}=\Gamma(\overline z)$ (conjugate the integral), so $|\Gamma(\frac12+iy)|^2=\Gamma(\frac12+iy)\Gamma(\frac12-iy)=\frac\pi{\sin(\frac\pi2+i\pi y)}=\frac\pi{\cos(i\pi y)}=\frac\pi{\cosh\pi y}$.
\end{proof}

So $\Gamma$ decays like $e^{-\pi|y|/2}$ along vertical lines: on the line $\operatorname{Re}z=\frac12$, exactly so. This decay is what makes
$\Gamma$ a useful factor in Chapter~10, where the zeta function is multiplied by $\pi^{-s/2}\Gamma(\frac s2)$. Every identity of Chapter~5 between functions
that extend holomorphically, the duplication formula for instance, extends to the complex plane by the same argument.
Exercise~\ref{ex:wielandt} gives the complex analogue of Bohr--Mollerup: $\Gamma$ is the only solution of $F(z+1)=zF(z)$, $F(1)=1$, that is
bounded on a vertical strip.

\section{Counting zeros}\label{sec:argument}

A function is \emph{meromorphic} on $U$ if it is holomorphic except at isolated points, all poles.

\begin{theorem}[argument principle]\label{thm:argument}
Let $D$ be a bounded convex open set with positive boundary $\gamma$, and $f$ meromorphic on an open set $U\supseteq\overline D$, with finitely many poles,
not identically zero on any disc, and with no zeros or poles on $\gamma^*$. Then
\[
\frac1{2\pi i}\int_\gamma\frac{f'(z)}{f(z)}\dd z=Z-P,
\]
where $Z$ and $P$ are the numbers of zeros and poles of $f$ in $D$, counted with multiplicity. The left side is also $n(f\circ\gamma,0)$.
\end{theorem}
\begin{proof}
We may shrink $U$ to the convex, hence connected, set $D_\varepsilon$ of the proof of Theorem~\ref{thm:residue}. The zeros of $f$ in the compact set $\overline D$ are
finitely many, since they are isolated (Theorem~\ref{thm:identity}, applied off the poles). At a zero of order $m$, $f=(z-a)^mg$ with $g(a)\ne0$, so
$\frac{f'}f=\frac m{z-a}+\frac{g'}g$ has residue $m$. At a pole of order $m$, $f=(z-a)^{-m}g$ with $g$ holomorphic and $g(a)\ne0$, and the residue is $-m$.
Elsewhere $\frac{f'}f$ is holomorphic. Apply Theorem~\ref{thm:residue}. Finally $\int_{f\circ\gamma}\frac{\mathrm dw}w=\int_\alpha^\beta\frac{f'(\gamma(t))\gamma'(t)}{f(\gamma(t))}\dd t$.
\end{proof}

So the number of zeros inside a contour is the number of times the image of the contour winds around $0$ (Figure~\ref{fig:winding}).
For a polynomial of degree $n$ and a huge circle, $f(z)\approx z^n$ and the image winds $n$ times: the fundamental theorem of
algebra again, now with the count. The next theorem makes ``$\approx$'' precise.

\begin{theorem}[Rouch\'e]\label{thm:rouche}
Let $D$, $\gamma$ be as above, $f$, $g$ holomorphic on an open set containing $\overline D$, and $|g|<|f|$ on $\gamma^*$. Then $f$ and $f+g$ have the same
number of zeros in $D$.
\end{theorem}
\begin{proof}
For $s\in[0,1]$, $f_s=f+sg$ satisfies $|f_s|\ge|f|-|g|>0$ on $\gamma^*$, so it is not identically zero, and the argument principle applies: its number of
zeros in $D$ is $N(s)=\frac1{2\pi i}\int_\gamma\frac{f'+sg'}{f+sg}$. The integrand is continuous in $(s,t)$ on the compact $[0,1]\times[\alpha,\beta]$ (piece by piece), and its
denominator is bounded below by $\min_{\gamma^*}(|f|-|g|)>0$; so $N$ is continuous by Theorem~\ref{thm:paramcompact}. An integer-valued continuous function on
$[0,1]$ is constant, and $N(0)=N(1)$.
\end{proof}

\begin{example}
How many roots of $z^5+3z+1$ lie in the unit disc? On $|z|=1$, $|z^5+1|\le2<3=|3z|$, so by Rouch\'e with $f=3z$ and $g=z^5+1$ there is exactly
one, like $3z$. On $|z|=\frac32$, $|3z+1|\le\frac{11}2<\frac{243}{32}=|z^5|$, so all five roots have modulus less than $\frac32$.
\end{example}

\begin{figure}[ht]
\centering
\fig{ch09-winding}
\caption{Left: the five zeros of $z^5+3z+1$ (red) and the circles $|z|=1$ (blue) and $|z|=\frac32$ (orange). Middle: the image of the unit circle
under $z\mapsto z^5+3z+1$, rescaled, winds once around $0$ (the black dot), despite its four loops. Right: the image of the larger circle winds
five times. Only one zero is inside the unit circle, and all five are inside the larger one.}
\label{fig:winding}
\end{figure}

\begin{think}
The argument principle turns a question about where the zeros are into an integral. Is this a practical algorithm for finding
roots? What would you do if the integral came out as $2.9999$? (A numerical calculation builds a root finder on this idea.)
\end{think}

\section{Exercises}

\begin{exercise}[$\star\star$]\label{ex:xn}
For $n\ge2$, show that $\displaystyle\int_0^\infty\frac{\mathrm dx}{1+x^n}=\frac{\pi/n}{\sin(\pi/n)}$, using the sector $0\le\arg z\le\frac{2\pi}n$, $|z|\le R$. Check the result against the
Beta function.
\end{exercise}

\begin{exercise}[$\star\star$]\label{ex:fresnelsector}
Integrating $e^{iz^2}$ over the sector $0\le\arg z\le\frac\pi4$, $|z|\le R$, show that $\displaystyle\int_0^\infty\cos(x^2)\dd x=\int_0^\infty\sin(x^2)\dd x=\frac12\sqrt{\frac\pi2}$. Compare with
Lemma~\ref{lem:fresnel}.
\end{exercise}

\begin{exercise}[$\star$]\label{ex:cothres}
Using Theorem~\ref{thm:residuesum}, show that $\displaystyle\sum_{n\ge1}\frac1{n^2+a^2}=\frac\pi{2a}\coth\pi a-\frac1{2a^2}$ for $a>0$. Compare with Exercise~\ref{ex:coth}. What is $\sum_{n\in\Z}\frac{(-1)^n}{n^2+a^2}$?
\end{exercise}

\begin{exercise}[$\star$]\label{ex:rouche}
Show that $z^4-6z+3$ has exactly one zero in the unit disc and three in the annulus $1<|z|<2$.
\end{exercise}

\begin{exercise}[$\star\star$]\label{ex:polygrowth}
Let $f$ be entire with $|f(z)|\le C(1+|z|)^k$ for all $z$. Show that $f$ is a polynomial of degree at most $k$. Deduce that an entire function $f$
with $|f(z)|\le C|e^z|$ for all $z$ is $ce^z$.
\end{exercise}

\begin{exercise}[$\star\star\star$]\label{ex:wielandt}
(Wielandt's theorem, 1939.) Let $F$ be holomorphic on $\operatorname{Re}z>0$ with $F(z+1)=zF(z)$, $F(1)=1$, and $F$ bounded on the strip $1\le\operatorname{Re}z\le2$. Show
that $F=\Gamma$.
\begin{enumerate}[label=(\alph*)]
\item Show that $\Gamma$ is bounded on the strip, and let $G=F-\Gamma$. Show that $G$ extends to an entire function satisfying $G(z+1)=zG(z)$.
\item Show that $H(z)=G(z)G(1-z)$ is bounded on $0\le\operatorname{Re}z\le1$ and satisfies $H(z+1)=-H(z)$. Conclude.
\end{enumerate}
\end{exercise}

\section{Hints}

\begin{hint}{ex:xn}
The only pole in the sector is $e^{i\pi/n}$. On the second ray $z=re^{2\pi i/n}$, $z^n=r^n$. For the check, substitute $u=x^n$ and use reflection.
\end{hint}
\begin{hint}{ex:fresnelsector}
On the arc, $|e^{iz^2}|=e^{-R^2\sin2\theta}$ and $\sin2\theta\ge\frac{4\theta}\pi$ on $[0,\frac\pi4]$ (concavity). On the ray $z=re^{i\pi/4}$, $iz^2=-r^2$.
\end{hint}
\begin{hint}{ex:cothres}
The residues of $\frac{\pi\cot\pi z}{z^2+a^2}$ at $\pm ia$; $\cot(iy)=-i\coth y$. For the alternating sum, $\frac\pi{\sin\pi z}$ has residue $(-1)^n$ at $n$.
\end{hint}
\begin{hint}{ex:rouche}
Compare with $-6z$ on $|z|=1$ and with $z^4$ on $|z|=2$. Check that there is no zero on $|z|=1$.
\end{hint}
\begin{hint}{ex:polygrowth}
Cauchy's estimates on $C(0,r)$, then $r\to\infty$. For the second part, look at $f(z)e^{-z}$.
\end{hint}
\begin{hint}{ex:wielandt}
(a) Log-convexity and $\Gamma(1)=\Gamma(2)=1$; near $-n$, write $G(z)=\frac{G(z+n+1)}{z(z+1)\cdots(z+n)}$ and use $G(1)=0$. (b) For $|\operatorname{Im}z|\ge1$ in the strip, $G(z)=\frac{G(z+1)}z$.
Then Liouville, and Theorem~\ref{thm:identity}.
\end{hint}

\section{Solutions}

\begin{solution}{ex:xn}
\idea The integrand is invariant under rotation by $\omega=e^{2\pi i/n}$, since $(\omega z)^n=z^n$: the integral along the rotated ray is $\omega$ times the
one we want, which is the same trick as the rectangle of Example~\ref{ex:rectangle}.
\writeup The sector is convex since its angle $\frac{2\pi}n$ is at most $\pi$. The zeros of $1+z^n$ are $e^{i\pi(2k+1)/n}$, and only $a=e^{i\pi/n}$ has argument in
$\left]0,\frac{2\pi}n\right[$; the residue there is $\frac1{na^{n-1}}=\frac a{na^n}=-\frac an$. On the arc, $\bigl|\frac1{1+z^n}\bigr|\le\frac1{R^n-1}$, and the arc has length $\frac{2\pi R}n$, so its integral
tends to $0$ as $n\ge2$. On the ray $z=r\omega$, traversed inwards, the integral is $-\omega\int_0^R\frac{\mathrm dr}{1+r^n}$. Letting $R\to\infty$, with $\omega=a^2$,
$(1-a^2)I=-\frac{2\pi ia}n$, so
\[
\begin{aligned}I&=\frac{2\pi i}{n(a-a^{-1})}\eqbr=\frac{2\pi i}{n\cdot2i\sin\frac\pi n}=\frac{\pi/n}{\sin(\pi/n)}.\end{aligned}
\]
Check: with $u=x^n$, $I=\frac1n\int_0^\infty\frac{u^{1/n-1}}{1+u}\dd u=\frac1n\Gamma(\frac1n)\Gamma(1-\frac1n)$ by Example~\ref{ex:rectangle}, and reflection gives the same value. For $n=2$, $\frac\pi2$.
\end{solution}

\begin{solution}{ex:fresnelsector}
\idea $e^{iz^2}$ oscillates on the real axis but, on the diagonal $z=re^{i\pi/4}$, it becomes the Gaussian $e^{-r^2}$. Cauchy's theorem lets us move from
one to the other, provided the arc between them contributes nothing.
\writeup $e^{iz^2}$ is entire, so its integral over the boundary of the sector (convex) is $0$. On the arc $z=Re^{i\theta}$, $iz^2=iR^2(\cos2\theta+i\sin2\theta)$, so
$|e^{iz^2}|=e^{-R^2\sin2\theta}\le e^{-4R^2\theta/\pi}$, using that the concave function $\sin2\theta$ lies above its chord on $[0,\frac\pi4]$. The arc integral is therefore at
most $\int_0^{\pi/4}e^{-4R^2\theta/\pi}R\dd\theta\le\frac{\pi R}{4R^2}=\frac\pi{4R}\to0$. On the ray, $z=re^{i\pi/4}$ and $iz^2=-r^2$, contributing $-e^{i\pi/4}\int_0^Re^{-r^2}\dd r$. So
\[
\int_0^Re^{ix^2}\dd x\longrightarrow e^{i\pi/4}\int_0^\infty e^{-r^2}\dd r=\frac{\sqrt\pi}2\cdot\frac{1+i}{\sqrt2},
\]
which proves that the improper integral converges, and taking real and imaginary parts gives $\frac{\sqrt\pi}{2\sqrt2}=\frac12\sqrt{\frac\pi2}$ for both. With $u=x^2$,
$\int_0^\infty e^{ix^2}\dd x=\frac12\int_0^\infty u^{-1/2}e^{iu}\dd u$, which is $\frac12\sqrt\pi e^{i\pi/4}$ by Lemma~\ref{lem:fresnel}: the same value, by a different road.
\end{solution}

\begin{solution}{ex:cothres}
\idea Apply the theorem with the two poles $\pm ia$ of $f$; the cotangent of an imaginary number is a hyperbolic cotangent.
\writeup $f(z)=\frac1{z^2+a^2}$ satisfies the bound and has simple poles at $\pm ia$, not integers. By Proposition~\ref{prop:rescalc}(b), the residue of
$\frac{\pi\cot\pi z}{z^2+a^2}$ at $\pm ia$ is $\frac{\pi\cot(\pm i\pi a)}{\pm2ia}$. Since $\cot(iy)=\frac{\cos iy}{\sin iy}=\frac{\cosh y}{i\sinh y}=-i\coth y$ and $\cot$ is odd, both residues equal
$\frac{-i\pi\coth\pi a}{2ia}=-\frac{\pi\coth\pi a}{2a}$. So $\sum_{n\in\Z}\frac1{n^2+a^2}=\frac\pi a\coth\pi a$, as in Exercise~\ref{ex:coth} by Poisson summation; removing $n=0$ and halving gives
the formula. For the alternating sum, $\frac\pi{\sin\pi z}$ also satisfies Lemma~\ref{lem:cotbound}'s kind of bound on $C_N$ (on the horizontal sides
$\bigl|\frac{2i}{e^{i\pi z}-e^{-i\pi z}}\bigr|\le\frac2{e^{\pi/2}-e^{-\pi/2}}$, on the vertical sides $|\sin\pi z|=\cosh\pi y\ge1$), and has residue $(-1)^n$ at $n$. The same proof gives
$\sum_n\frac{(-1)^n}{n^2+a^2}=-2\cdot\frac{\pi}{2ia\sin(i\pi a)}=\frac\pi{a\sinh\pi a}$, using $\sin(iy)=i\sinh y$.
\end{solution}

\begin{solution}{ex:rouche}
\idea Find on each circle the term that dominates the others.
\writeup On $|z|=1$, $|z^4+3|\le4<6=|-6z|$, so $z^4-6z+3$ has as many zeros in the unit disc as $-6z$, namely one; the strict inequality also shows
there is no zero on the circle. On $|z|=2$, $|-6z+3|\le15<16=|z^4|$, so there are four zeros in $|z|<2$, as for $z^4$. Hence three zeros in $1<|z|<2$.
\end{solution}

\begin{solution}{ex:polygrowth}
\idea Liouville's proof works with polynomial growth: Cauchy's estimates kill every coefficient beyond the degree.
\writeup $f(z)=\sum c_nz^n$ on $\C$ by Theorem~\ref{thm:analytic}, and Corollary~\ref{cor:radius}(a) gives $|c_n|\le\frac{C(1+r)^k}{r^n}$ for every $r>0$. For $n>k$ this tends to $0$ as
$r\to\infty$, so $c_n=0$. For the second part, $g(z)=f(z)e^{-z}$ is entire and bounded by $C$, hence constant by Liouville.
\end{solution}

\begin{solution}{ex:wielandt}
\idea On the real axis Bohr--Mollerup used log-convexity; in the complex plane boundedness on a strip takes its place, and
Liouville does the rest. The product $G(z)G(1-z)$ is built to be periodic, exactly as $\Gamma(z)\Gamma(1-z)$ is.
\writeup (a) $\Gamma$ is log-convex on $\left]0,\infty\right[$ (Corollary~\ref{cor:gammalogconvex}) with $\Gamma(1)=\Gamma(2)=1$, so $\Gamma\le1$ on $[1,2]$, and $|\Gamma(z)|\le\Gamma(\operatorname{Re}z)\le1$ on the
strip by Theorem~\ref{thm:gammac}(a). So $G$ is bounded on the strip, $G(z+1)=zG(z)$ and $G(1)=0$. As in Theorem~\ref{thm:gammac}(b), $G$ extends to
$\C\setminus\{0,-1,\dots\}$ by $G(z)=\frac{G(z+n+1)}{z(z+1)\cdots(z+n)}$. Near $-n$, the numerator tends to $G(1)=0$, so $(z+n)G(z)\to0$, and $G$ is bounded near $-n$ (the numerator is
holomorphic and vanishes at $-n$, so it is $(z+n)$ times a holomorphic function, by Theorem~\ref{thm:identity}(a)). By Theorem~\ref{thm:removable}, $G$ is entire.

(b) On the strip $S_0=\{0\le\operatorname{Re}z\le1\}$, $G(z)=\frac{G(z+1)}z$ is bounded for $|\operatorname{Im}z|\ge1$, and $G$ is bounded on the compact part $|\operatorname{Im}z|\le1$ by continuity.
As $z\mapsto1-z$ maps $S_0$ to itself, $H$ is bounded on $S_0$. Next, $G(1-z)=G(-z+1)=-zG(-z)$, so for $z\ne0$,
\[
\begin{aligned}H(z+1)&=G(z+1)G(-z)=zG(z)\cdot\frac{G(1-z)}{-z}\eqbr=-H(z),\end{aligned}
\]
and by continuity also at $z=0$. So $|H|$ is $1$-periodic, hence bounded on $\C$, and $H$ is constant by Liouville: $H=H(1)=G(1)G(0)=0$. If $G$ were not
identically zero, its zeros would be isolated, so $G\ne0$ on some disc $D$; then $G(1-z)=0$ for $z\in D$, and Theorem~\ref{thm:identity} gives $G=0$, a
contradiction. Hence $F=\Gamma$.
\end{solution}

\section*{Sources and further reading}

\begin{itemize}
\item E.~Stein and R.~Shakarchi, \emph{Complex Analysis}, Princeton University Press, 2003, Chapters~1 to~3 (Goursat,
Cauchy, residues, Rouch\'e) and~6 (the Gamma function).
\item L.~Ahlfors, \emph{Complex Analysis}, 3rd ed., McGraw-Hill, 1979. The classic, with the homology form of Cauchy's theorem.
\item T.~Needham, \emph{Visual Complex Analysis}, Oxford University Press, 1997. Pictures for everything in this chapter,
including winding numbers and Rouch\'e's ``dog on a leash''.
\item E.~Goursat, \emph{Sur la d\'efinition g\'en\'erale des fonctions analytiques, d'apr\`es Cauchy}, Transactions of the AMS 1
(1900), 14--16. Theorem~\ref{thm:goursat} without continuity of $f'$.
\item A.-L.~Cauchy, \emph{M\'emoire sur les int\'egrales d\'efinies prises entre des limites imaginaires}, Paris, 1825. The birth
of contour integration.
\item R.~Remmert, \emph{Wielandt's theorem about the $\Gamma$-function}, American Mathematical Monthly 103 (1996), 214--220.
Exercise~\ref{ex:wielandt}.
\item D.~S.~Mitrinovi\'c and J.~D.~Ke\v cki\'c, \emph{The Cauchy Method of Residues}, Reidel, 1984. Hundreds of integrals and
sums computed by the methods of this chapter.
\end{itemize}
